\chapter{สถาปัตยกรรมสมองกลฝังตัวและการควบคุมฮาร์ดแวร์กำลังสูงในแปลงเกษตร}
\label{chap:ch02}

\section*{แผนบริหารการสอนประจำบทที่ 2}
\addcontentsline{toc}{section}{แผนบริหารการสอนประจำบทที่ 2}

\begin{planbox}{หัวข้อเนื้อหาประจำบท}
\begin{enumerate}[topsep=2pt, itemsep=1pt, partopsep=0pt, parsep=0pt, leftmargin=1.5em]
    \item สถาปัตยกรรมหน่วยประมวลผลไมโครคอนโทรลเลอร์ ESP32-S3 สำหรับงานเกษตรกรรมแม่นยำ
    \item ระบบปฏิบัติการเวลาจริง FreeRTOS\index{ระบบปฏิบัติการเวลาจริง (FreeRTOS)}\index{FreeRTOS} การจัดลำดับความสำคัญของภารกิจ และระบบเฝ้าระวังวอทช์ด็อก
    \item วงจรขับและเทคนิคการควบคุมโซลินอยด์วาล์วไฟฟ้า 24VAC เพื่อป้องกันค้อนน้ำชลศาสตร์
    \item การควบคุมมอเตอร์ 3 เฟสด้วยแมกเนติกคอนแทกเตอร์ โอเวอร์โหลดรีเลย์ และสวิตช์แรงดัน
    \item การควบคุมความเร็วรอบมอเตอร์เครื่องตีน้ำและปั๊มน้ำด้วยอินเวอร์เตอร์ปรับความถี่ผ่าน Modbus RTU
\end{enumerate}
\end{planbox}

\begin{planbox}{วัตถุประสงค์เชิงพฤติกรรม}
\begin{enumerate}[topsep=2pt, itemsep=1pt, partopsep=0pt, parsep=0pt, leftmargin=1.5em]
    \item อธิบายสถาปัตยกรรมฮาร์ดแวร์ของชิป ESP32-S3 และการจัดสรรหน่วยความจำสำหรับงาน Edge AI ได้อย่างถูกต้อง
    \item ออกแบบวงจรขับโซลินอยด์วาล์ว 24VAC ด้วยออปโตคัปเปลอร์และไตรแอกได้อย่างปลอดภัย
    \item เขียนโปรแกรม FreeRTOS แบบหลายภารกิจเพื่อควบคุมการให้น้ำแบบไม่ปิดกั้นการทำงานของระบบ
    \item สื่อสารและกำหนดความถี่การหมุนของอินเวอร์เตอร์ VFD ผ่านรีจิสเตอร์ Modbus RTU ได้อย่างแม่นยำ
\end{enumerate}
\end{planbox}

\newpage

\begin{planbox}{วิธีสอนและกิจกรรมการเรียนการสอน}
\begin{enumerate}[topsep=2pt, itemsep=1pt, partopsep=0pt, parsep=0pt, leftmargin=1.5em]
    \item บรรยายเชิงทฤษฎีสถาปัตยกรรมไมโครโปรเซสเซอร์และวิศวกรรมไฟฟ้ากำลังควบคุมเครื่องกล
    \item สาธิตการต่อวงจรทดลองบอร์ดคอนโทรลเลอร์เข้ากับตู้ควบคุมไฟฟ้าอุตสาหกรรม
    \item ฝึกปฏิบัติการเขียนโค้ดภาษา MicroPython รันงานคู่ขนานแบบสองแกนประมวลผล
    \item ปฏิบัติการทดสอบควบคุมความถี่อินเวอร์เตอร์เพื่อปรับอัตราการหมุนของใบพัดตีน้ำ
\end{enumerate}
\end{planbox}

\begin{planbox}{สื่อการเรียนการสอน}
\begin{enumerate}[topsep=2pt, itemsep=1pt, partopsep=0pt, parsep=0pt, leftmargin=1.5em]
    \item เอกสารคำสอนบทที่ 2 ชุดสไลด์และผังวงจรไฟฟ้ากำลังควบคุม
    \item บอร์ดพัฒนา ESP32-S3 Dual-Core และชุดโมดูลรีเลย์อุตสาหกรรม Solid State
    \item อินเวอร์เตอร์ขับมอเตอร์ VFD ขนาด 2 แรงม้า พร้อมมอเตอร์เหนี่ยวนำ 3 เฟส
    \item โซลินอยด์วาล์วไฮดรอลิกส์ 24VAC ขนาด 2 นิ้ว พร้อมทรานส์ฟอร์เมอร์สเต็ปดาวน์
\end{enumerate}
\end{planbox}

\begin{planbox}{การวัดและประเมินผล}
\begin{enumerate}[topsep=2pt, itemsep=1pt, partopsep=0pt, parsep=0pt, leftmargin=1.5em]
    \item การทดสอบข้อเขียนวัดความรู้ความเข้าใจเรื่องระบบไฟฟ้ากำลังและความปลอดภัย
    \item การตรวจความถูกต้องของการต่อวงจรควบคุมและการต่อสายดินตามมาตรฐานสากล
    \item การประเมินสมรรถนะการเขียนโปรแกรมควบคุม VFD ผ่านคำสั่ง Modbus RTU
\end{enumerate}
\end{planbox}

\clearpage

\section{สถาปัตยกรรมระบบสมองกลฝังตัวประสิทธิภาพสูง}

ระบบสมองกลฝังตัวสำหรับการเกษตรอัจฉริยะและฟาร์มสัตว์น้ำยุคใหม่ จำเป็นต้องมีสมรรถนะการประมวลผลสัญญาณตัวเลขสูง สามารถรองรับการอนุมานโมเดลโครงข่ายประสาทเทียมขนาดเล็ก (TinyML) และมีระบบการสื่อสารไร้สายครบวงจร ชิปประมวลผล \textbf{ESP32-S3} ซึ่งใช้สถาปัตยกรรม 32-bit Xtensa LX7 สองแกนประมวลผล (Dual-core) ทำงานที่ความเร็วสัญญาณนาฬิกาสูงสุด $240\text{ MHz}$ พร้อมหน่วยคำนวณเวกเตอร์ (Vector Instructions) เหมาะอย่างยิ่งสำหรับการสกัดคุณลักษณะของภาพถ่ายและสัญญาณเสียง

\begin{table}[H]
\centering
\caption{คุณสมบัติทางวิศวกรรมของไมโครคอนโทรลเลอร์สำหรับการเกษตรอัจฉริยะ}
\label{tab:mcu_specs}
\begin{tabularx}{\textwidth}{p{3.5cm}YYY}
    \toprule
    \textbf{พารามิเตอร์} & \textbf{ESP32-S3} & \textbf{ESP32-C3} & \textbf{Raspberry Pi Pico W} \\
    \midrule
    สถาปัตยกรรมซีพียู & Xtensa LX7 (Dual-core) & RISC-V (Single-core) & ARM Cortex-M0+ (Dual) \\
    ความเร็วสัญญาณนาฬิกา & 240 MHz & 160 MHz & 133 MHz \\
    หน่วยประมวลผลเวกเตอร์ & มี (AI Vector Ext.) & ไม่มี & ไม่มี \\
    หน่วยความจำ SRAM & 512 KB + PSRAM สูงสุด 8MB & 400 KB & 264 KB \\
    การสื่อสารไร้สาย & Wi-Fi 4 + BLE 5.0 Mesh & Wi-Fi 4 + BLE 5.0 & Wi-Fi 4 + BLE 5.2 \\
    ความเหมาะสมในระบบ & Edge AI และวิทัศน์ & โหนดเซนเซอร์สนาม & โหนดควบคุมทั่วไป \\
    \bottomrule
\end{tabularx}
\end{table}

\section{ระบบปฏิบัติการเวลาจริง FreeRTOS และการจัดสรรภารกิจ}

ในระบบควบคุมฟาร์ม การอ่านค่าเซนเซอร์ การรับส่งข้อมูลผ่านโครงข่ายไร้สาย และการควบคุมวาล์วหรือมอเตอร์ ต้องทำงานพร้อมกันอย่างแม่นยำ การเขียนโปรแกรมแบบวนลูปเดี่ยวดั้งเดิมจะทำให้เกิดการปิดกั้นการทำงาน (Blocking) การนำระบบปฏิบัติการเวลาจริง \textbf{FreeRTOS} เข้ามาบริหารจัดการ ทำให้สามารถแยกงานออกเป็นภารกิจย่อย (Tasks) ที่มีลำดับความสำคัญต่างกันได้

\subsection{การจัดลำดับความสำคัญและระบบเฝ้าระวังวอทช์ด็อก}

ระบบควบคุมต้องมีการจัดลำดับความสำคัญของภารกิจ โดยกำหนดให้ภารกิจด้านความปลอดภัย เช่น การตรวจสอบแรงดันเกินในท่อน้ำ หรือการตรวจสอบระดับออกซิเจนวิกฤติ มีลำดับความสำคัญสูงสุด พร้อมเปิดใช้งาน \textbf{ระบบเฝ้าระวังวอทช์ด็อก (Task Watchdog Timer หรือ TWDT)} เพื่อรีเซ็ตระบบอัตโนมัติหากเกิดสภาวะซอฟต์แวร์ค้าง (Deadlock)

\begin{pythonbox}[การสร้างภารกิจ FreeRTOS แบบหลายแกนด้วย MicroPython (uasyncio)]
import uasyncio as asyncio
from machine import Pin, WDT

# กำหนดขาควบคุมฮาร์ดแวร์
VALVE_PIN = Pin(18, Pin.OUT)
PUMP_PIN = Pin(19, Pin.OUT)

async def telemetry_task():
    while True:
        # อ่านข้อมูลเซนเซอร์และส่งขึ้นคลาวด์ทุก 10 วินาที
        print("กำลังอ่านค่าเซนเซอร์และส่งข้อมูลเทเลเมทรี...")
        await asyncio.sleep(10)

async def valve_controller_task(wdt):
    while True:
        # ภารกิจควบคุมวาล์วและการป้อนวอทช์ด็อก
        wdt.feed()
        await asyncio.sleep(1)

async def main():
    wdt = WDT(timeout=8000) # กำหนดเวลาตัด 8 วินาที
    task1 = asyncio.create_task(telemetry_task())
    task2 = asyncio.create_task(valve_controller_task(wdt))
    await asyncio.gather(task1, task2)

# รันระบบภารกิจคู่ขนาน
# asyncio.run(main())
\end{pythonbox}

\section{การควบคุมโซลินอยด์วาล์ว 24VAC และการป้องกันค้อนน้ำ}

การควบคุมการให้น้ำในสวนผลไม้ขนาดใหญ่นิยมใช้ \textbf{โซลินอยด์วาล์วไฟฟ้าชนิด 24VAC} เนื่องจากมีความปลอดภัยต่อผู้ปฏิบัติงานในสวนที่มีความชื้นสูง และสายสัญญาณไม่เกิดปัญหาแรงดันตกคร่อมรุนแรงเหมือนระบบแรงดันต่ำกระแสตรง

การตัดต่อวาล์วน้ำขนาดใหญ่ (ตั้งแต่ 2 นิ้วขึ้นไป) อย่างกะทันหัน จะก่อให้เกิดการเปลี่ยนแปลงโมเมนตัมของมวลน้ำอย่างรวดเร็ว ส่งผลให้เกิดคลื่นแรงกระแทกไฮดรอลิกส์ หรือ \textbf{ปรากฏการณ์ค้อนน้ำ (Water Hammer)\index{ปรากฏการณ์ค้อนน้ำ (Water Hammer)}\index{Water Hammer}} ซึ่งสามารถสร้างแรงดันกระชากสูงกว่าแรงดันใช้งานปกติได้ถึง $3 - 5$ เท่า จนท่อและข้อต่อแตกเสียหาย

\begin{theorybox}[หลักการป้องกันปรากฏการณ์ค้อนน้ำในระบบวาล์วไฟฟ้า]
การแก้ไขปัญหาค้อนน้ำทำได้โดย:
\begin{enumerate}[leftmargin=1.5em]
    \item การเลือกใช้วาล์วไฟฟ้าชนิดปิดช้า (Slow-Closing Solenoid Valve) ที่มีห้องควบคุมไดอะแฟรมหน่วงเวลา 3 ถึง 5 วินาที
    \item การออกแบบอัลกอริทึมการเปิด-ปิดวาล์วแบบเหลื่อมเวลา (Staggered Valve Sequencing) โดยเปิดวาล์วโซนถัดไปก่อนปิดวาล์วโซนเดิม 5 วินาที เพื่อรักษาสมดุลแรงดันในเส้นท่อเมน
    \item การขับวาล์วด้วยออปโตคัปเปลอร์ชนิดตรวจจับแรงดันศูนย์ (Zero-Crossing Optoisolator เช่น MOC3041) ร่วมกับไตรแอก (TRIAC) เพื่อลดสัญญาณรบกวนแม่เหล็กไฟฟ้าขณะสวิตชิ่ง
\end{enumerate}
\end{theorybox}

\section{การควบคุมมอเตอร์ 3 เฟสด้วยอินเวอร์เตอร์ VFD ผ่าน Modbus}

ในฟาร์มเพาะเลี้ยงสัตว์น้ำ เครื่องตีน้ำใบพัดเป็นอุปกรณ์ที่ใช้พลังงานไฟฟ้าสูงสุด การใช้มอเตอร์เหนี่ยวนำ 3 เฟส ร่วมกับ \textbf{อินเวอร์เตอร์ปรับความถี่รอบ (Variable Frequency Drive หรือ VFD\index{อินเวอร์เตอร์ปรับความถี่รอบ (VFD)}\index{Variable Frequency Drive (VFD)})} ช่วยให้สามารถปรับความเร็วรอบของใบพัดตามความต้องการออกซิเจนจริงในแต่ละช่วงเวลาของวัน

ไมโครคอนโทรลเลอร์เชื่อมต่อไปยังพอร์ต RS485 ของอินเวอร์เตอร์ และส่งคำสั่งควบคุมความถี่การหมุนและทิศทางการหมุนผ่านรีจิสเตอร์ Modbus RTU ดังตัวอย่างในโค้ดต่อไปนี้

\begin{pythonbox}[ฟังก์ชันปรับความเร็วรอบอินเวอร์เตอร์ VFD ด้วยภาษา Python]
import struct

def build_vfd_write_frequency_packet(slave_id: int, target_hz: float) -> bytes:
    # ค่าความถี่คูณด้วย 100 เพื่อแปลงเป็นจำนวนเต็ม เช่น 50.00 Hz = 5000
    freq_value = int(target_hz * 100)
    # ฟังก์ชัน 0x06 เขียนรีจิสเตอร์เดี่ยว รีจิสเตอร์ 0x2001 (ความถี่เป้าหมาย)
    packet = struct.pack(">BBHH", slave_id, 0x06, 0x2001, freq_value)
    # คำนวณรหัส CRC-16 Checksum
    crc = calculate_crc16(packet)
    packet += struct.pack("<H", crc)
    return packet

# ตัวอย่างการใช้งาน สั่งปรับความถี่ที่ตัวแปรอินเวอร์เตอร์หมายเลข 1 เป็น 42.50 Hz
# cmd = build_vfd_write_frequency_packet(1, 42.50)
\end{pythonbox}

\section*{คำถามทบทวนท้ายบทที่ 2}
\addcontentsline{toc}{section}{คำถามทบทวนท้ายบทที่ 2}

\begin{enumerate}
    \item จงอธิบายบทบาทของหน่วยประมวลผลเวกเตอร์ในชิป ESP32-S3 ว่ามีส่วนช่วยเร่งความเร็วในการประมวลผลปัญญาประดิษฐ์ที่ขอบอย่างไร
    \item ปรากฏการณ์ค้อนน้ำ (Water Hammer) เกิดจากสาเหตุใด และสามารถใช้หลักการควบคุมเชิงซอฟต์แวร์ป้องกันได้อย่างไร
    \item เหตุใดการขับโซลินอยด์วาล์ว 24VAC จึงต้องใช้โซลิดสเตตรีเลย์หรือไตรแอกชนิด Zero-Crossing
    \item หากต้องการลดการใช้พลังงานไฟฟ้าของเครื่องตีน้ำบ่อกุ้งลงร้อยละ 40 ควรปรับลดความถี่ของอินเวอร์เตอร์ VFD ลงเท่าใดตามกฎของพัดลมและปั๊ม (Affinity Laws)
\end{enumerate}
\clearpage
